\documentclass[10pt,conference,compsocconf]{IEEEtran}



\usepackage[dvipsnames]{xcolor}
\usepackage[utf8]{inputenc}   % кодировка исходного файла
\usepackage[T2A]{fontenc}     % русские буквы в PDF
\usepackage[russian]{babel}   % переносы, названия "Рис." и т.п.

\PassOptionsToPackage{hyphens}{url}\usepackage[
colorlinks = true,
linkcolor = BlueViolet,
urlcolor  = BlueViolet,
citecolor = BlueViolet,
anchorcolor = BlueViolet
]{hyperref}
\usepackage{graphicx}	% For figure environment
\usepackage{booktabs}
\usepackage{multirow}
\usepackage{tabularx}
\usepackage{amsmath}
\usepackage{amsfonts}


\begin{document}
\title{Двухканальное разделение речи на временных сигналах с помощью Dual-Path RNN
}

\author{
  М. Афанасьев, Д. Барсуков, А. Ковыляев\\
  \textit{Faculty Of Computer Science, HSE, Moscow}\\ 
}

\maketitle

\begin{abstract}
Направление разделения источников звука из единого канала на отдельные аудио-дорожки на сегодняшний день активно развивается во множестве областей. В данной работе рассматривается задача разделения голосов двух дикторов, на основе записанного аудио микса. Целью исследования было продвинуться в качестве решения данной задачи за счет изучения и применения архитектуры, учитывающей временной параметр в своих предсказаниях. Для этой цели мы взяли рекуррентную нейронную сеть с двойным путём (Dual-Path Recurrent Neural Networks, DPRNN)~\cite{luo_dprnn} в которую передаются закодированые аудио-последовательности, обработанные сетью разделения аудиосигналов во временной области (Time-domain Audio Separation Network, TasNet)~\cite{luo_tasnet} и протестировали её на новом наборе данных из нашей задачи. Качество модели оценивалось с точки зрения зашумленности получившихся аудио-последовательностей, на ряду с размерами и скоростью обработки входных данных. 

\end{abstract}

\section{Вступление}

Разделение источников звука на моно сигналы - это задача, которая появляется в различных практических приложениях, таких как видеоконференции, шумоподавление и других. В её основе нам задается моносигнал, содержащий в себе записи нескольких источников звука. На выходе мы ожидаем увидеть несколько каналов, каждый из которых содержит в себе отдельную запись каждого самостоятельного источника, которая не зашумлена другими сигналами. За последние годы инструменты, для решения данной проблемы значительно продвинулись благодаря развитию глубинного обучения. 
Долгое время доминировали методы, работающие в частотной области и использующие спектрограммы (например, на основе кратковременного преобразования Фурье (Short-Time Fourier Transform, STFT)~\cite{allen_stft} и маскирование амплитудных спектров. Однако появление полностью временных архитектур показало, что прямое считывание сигнала во времени позволяет избежать потерь информации и ограничений, связанных с обратимыми преобразованиями. Одним из наиболее успешных временных подходов стал DPRNN, где длинная последовательность разбивается на блоки, а временные зависимости моделируются отдельно в пределах каждого блока и между блоками. В своей работе мы использовали описанную выше архитектуру, для решения проблемы разделения двух источников речи, слитых в общий звуковой ряд. 
	Для сравнения и валидации подходов использовалась метрика Масштабно-инвариантное отношение сигнал/шум (Scale-Invariant Signal-to-Noise Ratio, Si-SNR), отражающая чистоту полученного сигнала. В отличие от классического SNR эта метрика не зависит от масштаба амлитуды. 
\section{Постановка задачи и данные}

Целью работы является исследование и реализация нейросетевой модели
одно\-канального разделения речи, способной по входному аудиосигналу
$x(t)$, содержащему сумму двух голосов,
\[
x(t) = s_1(t) + s_2(t),
\]
восстановить оценки исходных сигналов $\hat{s}_1(t)$ и $\hat{s}_2(t)$,
максимально близкие к истинным $s_1(t)$ и $s_2(t)$ в смысле метрики
SI-SNR и её приращения SI-SNRi.

Для обучения и валидации используется кастомный датасет, содержащий:
\begin{itemize}
  \item смешанные сигналы (папка \texttt{mix});
  \item отдельные каналы для первого и второго говорящего (папки \texttt{s1} и \texttt{s2});
  \item сопутствующую структуру разбиения на обучающую и валидационную части.
\end{itemize}

\section{Методология}

\subsection{Общий обзор архитектуры DPRNN}

DPRNN работает во временной области и строится вокруг идеи разбиения
латентного представления аудиосигнала на перекрывающиеся сегменты фиксированной длины.
После прохождения через энкодер исходная смесь $x(t)$
преобразуется в последовательность признаков
\[
\mathbf{Y} = [\mathbf{y}_1, \dots, \mathbf{y}_T] \in \mathbb{R}^{T \times C},
\]
где $C$ — размерность скрытого пространства. Далее последовательность
разбивается на чанки длиной $K$ с перекрытием, в результате чего
образуется тензор размера $(B, N, K, C)$, где $N$ — число сегментов.

Каждый блок DPRNN состоит из двух последовательно применяемых
рекуррентных подсетей:
\begin{itemize}
  \item \emph{внутрисегментной} (intra-chunk) RNN, которая моделирует
  зависимости внутри каждого сегмента длиной $K$;
  \item \emph{межсегментной} (inter-chunk) RNN, которая моделирует
  зависимости между сегментами по измерению $N$.
\end{itemize}
Обе RNN, как правило, реализуются двунаправленными LSTM с остаточными
соединениями и нормализацией по каналам, что обеспечивает стабильность
обучения и хорошую сходимость.

На выходе цепочки DPRNN-блоков полученное скрытое представление
подаётся в маскирующий модуль, который для каждого из двух источников
генерирует маску над латентными признаками. Применяя эти маски
к энкодерному представлению, мы получаем две раздельные последовательности,
которые затем декодируются обратно во временную область.

\subsection{Реализация модели и кода}

Наша реализация DPRNN в виде
модульной архитектуры:
\begin{itemize}
  \item \textbf{Код модели.}
  DPRNN разделена на независимые классы: энкодер, DPRNN-блоки,
  маскирующий модуль и декодер. Это упрощает модификации
  отдельных компонентов (например, изменение числа блоков или
  размера латентного пространства) без переписывания всей модели.
  \item \textbf{Конфигурирование через Hydra.}
  Все ключевые гиперпараметры (длина чанка, число DPRNN-блоков,
  размер скрытого слоя, выбор функции потерь, параметры оптимизатора)
  задаются в YAML-конфигурациях и передаются в код через Hydra.
  Это позволяет воспроизводить эксперименты и систематически
  исследовать влияние различных настроек.
  \item \textbf{Тренировочный цикл.}
  Обучение организовано через базовый класс \texttt{BaseTrainer},
  который реализует общую логику эпох, логирования, сохранения чекпоинтов
  и ранней остановки. Специализированный \texttt{Trainer} реализует
  конкретные шаги для задачи разделения речи: вызов модели, вычисление
  функции потерь SI-SNR (с PIT или без) и расчёт метрик на батче.
\end{itemize}

\subsection{Функция потерь и метрики}

В качестве основной функции потерь используется SI-SNR, реализованная
через модуль \texttt{PermutationInvariantTraining}
из \texttt{torchmetrics}, который автоматически выбирает оптимальное
сопоставление предсказанных и истинных источников и максимизирует
SI-SNR по всем возможным перестановкам.


Качество работы модели оценивается по SI-SNRi и SDRi, вычисляемым
как разность между значениями SI-SNR/SDR для выходных сигналов
и для исходной смеси. Для анализа воспринимаемого качества и разборчивости
речи могут также использоваться PESQ и STOI~\cite{itu_pesq,taal_stoi}, но в данном отчёте
основной акцент сделан на SI-SNRi и SDRi.

\section{Эксперименты}

\subsection{Общая схема экспериментов}

Все модели обучались на одном и том же датасете аудиоразделения (двухголосые смеси речи). Для обеспечения сопоставимости результатов мы использовали единый пайплайн: одинаковый лосс (SI-SNR с permutation invariant training), одинаковое разбиение на обучающую и валидационную части, а также общий набор метрик качества (SI-SNRi, SDRi, PESQ, STOI).  

Основной целью серии экспериментов было последовательно улучшить качество DPRNN-модели за счёт изменения архитектуры и гиперпараметров а также сравнить DPRNN с сверточной архитектурой ConvTasNet и RTFS-Net~\cite{luo_convtasnet,pegg_rtfsnet} в идентичных условиях обучения.

\subsection{Эксперименты с DPRNN}

\subsubsection{Базовая DPRNN}

В качестве отправной точки мы взяли DPRNN-модель с архитектурой, близкой к описанной в оригинальной статье.  

Для базовой конфигурации использовались умеренные размеры:
\begin{itemize}
    \item размер эмбеддинга энкодера: 512
    \item бутылочное сужение DPRNN: 128
    \item число DPRNN-блоков: 4
    \item стандартный шаг обучения из базового конфига (без дополнительного тюнинга)
\end{itemize}

После обучения на ограниченном числе эпох модель достигла примерно \textbf{SI-SNRi $\approx 2$~} на валидации. Кривая лосса убывала, но довольно быстро выходила на плато, а качество оставалось существенно ниже ожидаемого уровня.

\begin{figure}[htbp]
    \centering
    \includegraphics[width=0.9\linewidth]{dprnn_1.jpeg}
    \caption{Динамика метрик для базовой DPRNN.}
    \label{fig:conv-metrics}
\end{figure}

\subsubsection{Изменение градиентного клиппинга}

Далее мы попробовали стабилизировать за счёт менее агрессивного градиентного клиппинга:
\begin{itemize}
    \item было выставлено ограничение \verb|max_grad_norm = 5|.
\end{itemize}

На практике это привело к обратному эффекту:
\begin{itemize}
    \item лосс уменьшался медленнее
    \item итоговое SI-SNRi оказалось даже ниже, чем у базовой модели без клиппинга
\end{itemize}

Таким образом, для нашей реализации DPRNN ограничение нормы градиента ухудшило оптимизацию, и мы отказались от этого варианта в дальнейшем.

\begin{figure}[htbp]
    \centering
    \includegraphics[width=0.9\linewidth]{dprnn_2.jpeg}
    \caption{Динамика метрик для DPRNN с grad norm 5.}
    \label{fig:conv-metrics}
\end{figure}

\subsubsection{Увеличение мощности модели}

Следующим этапом мы увеличили размеры DPRNN: мы увеличили число блоков и/или размер скрытых слоёв. Все остальные гиперпараметры оставались близкими к базовым.

Результат:
\begin{itemize}
    \item качество действительно немного выросло по сравнению с первой конфигурацией DPRNN
    \item SI-SNRi стал заметно выше 2, но всё ещё низким в контексте целевых значений для задачи разделения речи
    \item на валидации наблюдалась тенденция к переобучению: улучшения на трейне не полностью переносились на валидацию
\end{itemize}

Этот эксперимент показал, что увеличение числа параметров само по себе не решает проблему; ключевым становится подбор режима оптимизации.

\begin{figure}[htbp]
    \centering
    \includegraphics[width=0.9\linewidth]{dprnn_3.jpeg}
    \caption{Динамика метрик для DPRNN увеличенная.}
    \label{fig:conv-metrics}
\end{figure}

\subsubsection{Тюнинг learning rate и финальная DPRNN}

Понимая, что большая модель обладает достаточной предсказывающей способностью, мы сфокусировались на настройке шага обучения. На той же увеличенной DPRNN были проведены эксперименты с уменьшением \emph{learning rate} (с сохранением остального пайплайна: SI-SNR+UPIT, те же данные, без агрессивного клиппинга).

Обнаружив удачную комбинацию размера модели и learning rate, мы запустили финальное обучение на \textbf{200 эпох}. В этой конфигурации модель достигла \textbf{SI-SNRi $> 6$~} на валидационном наборе. Эту версию DPRNN мы рассматриваем как наш основной референс среди рекуррентных архитектур.

\begin{figure}[htbp]
    \centering
    \includegraphics[width=0.9\linewidth]{dprnn_final.jpeg}
    \caption{Динамика метрик для DPRNN увеличенной + lr.}
    \label{fig:conv-metrics}
\end{figure}

\subsection{Эксперименты с ConvTasNet}

В качестве сверточного базлайна мы реализовали и интегрировали модифицированную ConvTasNet-подобную модель. Архитектура строится по тому же принципу «энкодер–сепаратор–декодер», что и DPRNN, но вместо рекуррентных блоков использует стек TCN-блоков с дилатированными depthwise-свёртками и skip-коннектами.  

Важные моменты:
\begin{itemize}
    \item мы использовали те же TasNet-энкодер и декодер, что и в DPRNN, чтобы различия шли именно за счёт сепаратора;
    \item лосс и метрики полностью совпадают с DPRNN-экспериментами (SI-SNR с UPIT, SI-SNRi/SDRi/PESQ/STOI на валидации);
    \item режим обучения и датасет идентичны финальной конфигурации DPRNN.
\end{itemize}

На текущем этапе ConvTasNet показывает качество \emph{как минимум не хуже, а по ряду прогонов — даже лучше}, чем наша финальная версия DPRNN (SI-SNRi $> 6$~), при том что архитектура остаётся сравнительно простой и удобной для дальнейшей оптимизации. Дообучая и подбирая параметры модели мы в дальнейшем сможем достигнуть на данной архитектуре даже более высоких результатов.

\begin{figure}[htbp]
    \centering
    \includegraphics[width=0.9\linewidth]{convtasnet.jpeg}
    \caption{Динамика метрик для ConvTasNet.}
    \label{fig:conv-metrics}
\end{figure}

\subsection{Попытка воспроизведения RTFS-Net}

Первым шагом мы интегрировали в репозиторий модифицированный вариант RTFS-Net, сочетающий частотно-временную LSTM-обработку и визуальные признаки (FiLM-fusion для двух лиц). Модель была подключена к существующему трейнеру, лоссу SI-SNR и метрикам.

Несмотря на попытки стабилизировать обучение (ограничение длины аудио, подбор learning rate, упрощение архитектуры), RTFS-Net в нашей конфигурации так и не сошёлся до приемлемого качества. В дальнейшем мы рассматривали его как отрицательный результат и сосредоточились на DPRNN и ConvTasNet.

\section{Выводы и дальнейшая работа}
В данном отчёте мы рассмотрели постановку задачи одно\-канального
разделения речи и описали архитектуру DPRNN, ориентированную на
работу во временной области с учётом как локальных, так и глобальных
зависимостей в аудиосигнале. В результате удалось достигнуть хорошего качества на нашей задаче. Полный список результатов можно найти в Приложении, таблица 1.
В дальнейшем мы также планируем более детально протестировать архитектуру RTFS-Net, которая будет использовать дополнительную информацию - видеозаписи движения губ каждого спикера. Дополнительный слой параметров может значительно улучшить качество предсказаний модели за счет более глубокого понимания источников сигнала.

\section{Contribution}
\begin{itemize}
    \item М. Афанасьев: RTFS-Net, ConvTasNet, Dataset, Report
    \item Д. Барсуков: Обучение моделей, подбор гиперпараметров, тестирование моделей
    \item А. Ковыляев: DPRNN, Dataset, Demo, Inference
\end{itemize}
\begin{thebibliography}{99}

\bibitem{luo_dprnn}
Y.~Luo, C.~Han, N.~Mesgarani,
``Dual-Path RNN: Efficient Long Sequence Modeling for Time-Domain Single-Channel Speech Separation,''
in \textit{Proc. ICASSP}, 2020.

\bibitem{luo_tasnet}
Y.~Luo, N.~Mesgarani,
``TasNet: Time-Domain Audio Separation Network for Real-Time, Single-Channel Speech Separation,''
in \textit{Proc. ICASSP}, 2018.

\bibitem{pegg_rtfsnet}
D.~Pegg, X.~Li, B.~Hu,
``RTFS-Net: Recurrent Time-Frequency Modelling for Efficient Audio-Visual Speech Separation,''
in \textit{Proc. ICLR}, 2024.

\bibitem{itu_pesq}
ITU-T Recommendation P.862,
``Perceptual Evaluation of Speech Quality (PESQ): An Objective Method for End-to-End Speech Quality Assessment of Narrow-band Telephone Networks and Speech Codecs,''
International Telecommunication Union, 2001.

\bibitem{taal_stoi}
C.~H.~Taal, R.~C.~Hendriks, R.~Heusdens, J.~Jensen,
``A Short-Time Objective Intelligibility Measure for Time-Frequency Weighted Noisy Speech,''
in \textit{Proc. ICASSP}, 2011.

\bibitem{kingma_adam}
D.~P.~Kingma, J.~Ba,
``Adam: A Method for Stochastic Optimization,''
arXiv preprint arXiv:1412.6980, 2014.

\bibitem{ephrat_looking_to_listen}
A.~Ephrat \textit{et al.},
``Looking to Listen at the Cocktail Party: A Speaker-Independent Audio-Visual Model for Speech Separation,''
\textit{ACM Transactions on Graphics (SIGGRAPH)}, vol.~37, no.~4, 2018.

\bibitem{luo_convtasnet}
Y.~Luo, N.~Mesgarani,
``Conv-TasNet: Surpassing Ideal Time-Frequency Magnitude Masking for Speech Separation,''
\textit{IEEE/ACM Transactions on Audio, Speech, and Language Processing}, vol.~27, no.~8, pp.~1256--1266, 2019.

\bibitem{allen_stft}
J.~B.~Allen,
``Short Term Spectral Analysis, Synthesis, and Modification by Discrete Fourier Transform,''
\textit{IEEE Transactions on Acoustics, Speech, and Signal Processing}, vol.~25, no.~3, pp.~235--238, 1977.

\end{thebibliography}

\clearpage
\appendix
\section{Приложение A. Дополнительные эксперименты}
\label{sec:appendix-experiments}
\begin{table}[htbp]
    \centering
    \caption{Сравнение протестированных моделей по качеству разделения речи. 
    Значения приведены на валидационном наборе;}
    \label{tab:models-results}
    \begin{tabular}{lcccc}
        \toprule
        Модель & SI-SNRi $\uparrow$ & SDRi $\uparrow$ & PESQ $\uparrow$ & STOI $\uparrow$ \\
        \midrule
        DPRNN (базовая) 
            & $\approx 2.0$      & $\approx 2.3$     & 1.22 & 0.68 \\
        DPRNN + grad clip ($\text{max\_grad\_norm}=5$) 
            & $< 1.0$            & $< 1.0$           & 1.19 & 0.65 \\
        DPRNN (увеличенная модель) 
            & $\approx 1.2$      & $\approx 1.21$     & 1.29 & 0.66 \\
        DPRNN (увеличенная, уменьшенный LR, 200 эпох) 
            & $\mathbf{6.0+}$    & $\mathbf{5.5+}$   & $\mathbf{1.42}$ & $\mathbf{0.80}$ \\
        ConvTasNet (наш вариант) 
            & $\approx 2.9$    & -   & 1.24& 0.72 \\
        RTFS-Net (аудио-визуальная) 
            & -        & -      & - & - \\
        \bottomrule
    \end{tabular}
\end{table}


\end{document}
